% =============================================================================
% §8 LIMITATIONS + CONCLUSION — page budget: 0.5 pages (paper_plan_v2 §9)
% Written SIXTH, before the abstract (§16 P11 writing order).
%
% Content sources (plan §9 row 8 — all three named honestly):
%   * domains where per-step quality is intrinsically EXPENSIVE (SWE-bench: a
%     test-suite run per step; intermediate patches score ~0 until late,
%     degenerating τ* — plan §2.6/§5.1): the label machinery does not transfer
%     there; inference-time monitor on a frozen 32B agent is appendix-only.
%   * serving-regime dependence of the overhead story (KV-cache-fork vs
%     black-box re-prefill can FLIP the sign of stopper overhead — plan §5.3,
%     LearnStop 2606.30852): recommendation is regime-conditional (H6).
%   * scooping-window honesty (plan §8 risks: DASH/RePro/SlimSearcher/BAGEN are
%     weeks old; CMU CaRT/MRT group adjacent) — state what was concurrent.
%   * conclusion: restate the one claim (plan §0) + the honest boundary of what
%     internalization was shown to do (H5 verdict), NOT a general-AI flourish.
% =============================================================================
\section{Limitations and Conclusion}
\label{sec:limitations}

\TODO{Limitation 1: per-step-quality-expensive domains (SWE) --- why the label
machinery stops at the boundary (plan \S2.6).}

\TODO{Limitation 2: serving-regime dependence of overhead; regime-conditional
recommendation per H6 (plan \S5.3).}

\TODO{Limitation 3: concurrent-work window disclosed (plan \S8).}

\TODO{Conclusion: the bridge claim as proven (or its fallback framing as
realized); internalization boundary per the H5 verdict.}
